\documentclass[10pt,a4paper]{amsart}
\usepackage[margin=19mm]{geometry}
\usepackage{amsmath,amssymb,mathtools}
\usepackage[scr=rsfs,cal=euler]{mathalfa}
\usepackage{xfrac}
\usepackage[unicode,hidelinks,bookmarks=false]{hyperref}
\newcommand{\ie}{i.e.\ }
\newcommand{\cf}{cf.\ }
\newcommand{\resp}{respectively\ }
\newcommand{\Subsection}{\subsection}
\providecommand{\nobreakdash}{\nobreak\hbox{-}}
\setlength{\emergencystretch}{2em}
\multlinegap0pt
\usepackage{bm}
\usepackage{bbm}
\usepackage{dsfont}
\usepackage{xspace}
\usepackage{cancel}
\usepackage{mathtools}
\usepackage{amsthm}
\makeatletter
\@ifundefined{theorem}{\newtheorem{theorem}{Theorem}}{}
\makeatother
\title[Velocity Reconstruction from Flow-Induced Magnetic Fields]{Velocity Reconstruction from Flow-Induced Magnetic Fields}
\author{Yacine Mokhtari, Christina Frederick, Yunan Yang, Bjorn Engquist}
\date{Source: arXiv:2602.22097v2 (CC BY 4.0)}
\begin{document}
\maketitle

\noindent\emph{Source and licence.} Velocity Reconstruction from Flow-Induced Magnetic Fields. Version arXiv:2602.22097v2 (CC BY 4.0), distributed under the Creative Commons Attribution 4.0 International licence, as recorded in the arXiv licence metadata for this exact version and shown on the abstract page https://arxiv.org/abs/2602.22097v2. Full rights notice: licenses/arxiv-2602.22097-CC-BY-4.0.txt.\par\smallskip

\noindent\textit{Editorial scope.} Counted unit: the Theorem whose source label is \texttt{thm:exist-unique-Rd} in the article (source0.tex lines 347--389, source label \texttt{thm:exist-unique-Rd}), delivered together with the article's own complete original proof of it (source0.tex lines 391--450, one \texttt{proof} environment of 60 lines). No mathematics is written, repaired or re-derived: statement and proof are the article's own text line for line. Required context retained uncounted: none. Every definition, display and label of those lines is retained unchanged. Omitted from this package and not redistributed: 1--346, 390--390, 451--1164. Nothing inside a retained excerpt is omitted or altered: every retained body is delivered exactly as the article's lines are written, so no omission receipt of any kind accompanies this package. Citations: the retained excerpts carry no citation command at all, so no bibliography entry of the article is transcribed and no reference is redistributed. Reference adaptation: none was needed and none was made; every reference inside the delivered unit resolves inside it, so no unresolved reference or citation is produced. Printed slips kept literal: no printed word, symbol, hypothesis or number of the counted statement or of its proof was altered. Layout and class: the article's own class is \texttt{article} with packages including graphicx, bbm, amsmath, amssymb, amsthm, geometry, xcolor, hyperref; the wrapper is the lane's pinned \texttt{amsart} editorial wrapper at 10pt on A4, which restates the theorem-like environments the retained text needs and adds the article's own macro definitions that the retained text needs (none), with extra wrapper packages (bm, bbm, dsfont, xspace, cancel, mathtools). The delivered numbering is the unit's own. Assets: no retained span contains a figure environment, a graphics-inclusion command, a PDF-page inclusion, a table or a TikZ picture; no image file is copied into this package, the package declares no asset dependency and no image file is redistributed with it. Licence: version arXiv:2602.22097v2 of this article is distributed under the Creative Commons Attribution 4.0 International licence, as recorded in the arXiv licence metadata for this exact version and shown on the abstract page https://arxiv.org/abs/2602.22097v2. A full rights notice is delivered at licenses/arxiv-2602.22097-CC-BY-4.0.txt. Characters: every retained excerpt is pure ASCII, so the delivered engine, XeLaTeX with its default Latin Modern text font, reports no missing character.
\begin{theorem}[Existence and uniqueness on $\mathbb{R}^d$]
\label{thm:exist-unique-Rd}
Let $S\subset\mathbb R^d$ be a smooth $(d-1)$-dimensional hypersurface with unit normal
$\boldsymbol n(\boldsymbol x)$ such that
\[
\mathbf F\cdot \boldsymbol n(\boldsymbol x)\neq 0
\qquad \forall \boldsymbol x\in S,
\]
and assume that $S$ is globally transverse to the flow generated by $\mathbf F$, in the sense that
\[
\Phi:S\times\mathbb R\to\mathbb R^d,\qquad \Phi(y,s)=y+s\mathbf F,
\]
is a $C^1$-diffeomorphism. For $\boldsymbol x\in\mathbb R^d$, write uniquely
\[
\boldsymbol x=\pi(\boldsymbol x)+\tau(\boldsymbol x)\mathbf F,\qquad \pi(\boldsymbol x)\in S,\ \tau(\boldsymbol x)\in\mathbb R.
\]
Let $\boldsymbol h\in L^2(\mathbb R^d;\mathbb R^d)$.

\begin{enumerate}
\item[(i)]
There exists a function $\boldsymbol v\in L^2_{\mathrm{loc}}(\mathbb R^d;\mathbb R^d)$ such that {{}
\[
(\mathbf F\cdot\nabla)\boldsymbol v=\boldsymbol h\,.
\]}
Moreover, for any $\boldsymbol v_S\in L^2(S;\mathbb R^d)$, the function
\[
\boldsymbol v(\boldsymbol x)=\boldsymbol v_S(\pi(\boldsymbol x))+\int_{0}^{\tau(\boldsymbol x)}\boldsymbol h(\pi(\boldsymbol x)+s\mathbf F)\,ds
\]
defines a distributional solution with trace $\boldsymbol v|_S=\boldsymbol v_S$.

\item[(ii)]
For each prescribed trace $\boldsymbol v_S\in L^2(S;\mathbb R^d)$, there is at most one
$\boldsymbol v\in L^2_{\mathrm{loc}}(\mathbb R^d;\mathbb R^d)$ such that
\[
(\mathbf F\cdot\nabla)\boldsymbol v={{} \boldsymbol h},
\qquad
\boldsymbol v|_S=\boldsymbol v_S.
\]
\item[(iii)]
If in addition $\nabla\cdot\boldsymbol h=0$ on $\mathbb{R}^d$ and $\nabla\cdot\boldsymbol v= 0$ on $S$,
then we have $\nabla\cdot\boldsymbol v=0$ on $\mathbb{R}^d$.
\end{enumerate}
\end{theorem}

\begin{proof}
By assumption, $\Phi:S\times\mathbb R\to\mathbb R^d$, $\Phi(y,t)=y+t\mathbf F$, is a $C^1$-diffeomorphism.
Let $\Phi^{-1}(\boldsymbol x)=(\pi(\boldsymbol x),\tau(\boldsymbol x))$, so that every $\boldsymbol x\in\mathbb R^d$ admits the unique representation
$\boldsymbol x=\pi(\boldsymbol x)+\tau(\boldsymbol x)\mathbf F$ with $\pi(\boldsymbol x)\in S$ and $\tau(\boldsymbol x)\in\mathbb R$.

We first establish part \textit{(i)}. Define $\boldsymbol v$ by choosing any trace datum, for instance
$\boldsymbol v_S\equiv 0$, and setting
\begin{equation}\label{eq:v_def_Rd}
\boldsymbol v(\boldsymbol x):={{} \int_0^{\tau(\boldsymbol x)} \boldsymbol h(\pi(\boldsymbol x)+s\mathbf F)\,ds.}
\end{equation}
For each fixed $y\in S$, consider the function $\boldsymbol w_y:\mathbb R\to\mathbb R^d$ given by
$\boldsymbol w_y(s):=\boldsymbol v(y+s\mathbf F)$. By \eqref{eq:v_def_Rd},
\[
\boldsymbol w_y(s)={{} \int_0^s \boldsymbol h(y+s'\mathbf F)\,ds',}
\]
so $\boldsymbol w_y$ is absolutely continuous and satisfies
\[
\frac{d}{ds}\boldsymbol w_y(s)={{} \boldsymbol h(y+s\mathbf F)\,,}
\]
for a.e.~$s\in\mathbb{R}$. Since $\mathbf F$ is constant, the chain rule gives $\frac{d}{ds}\boldsymbol v(y+s\mathbf F)=(\mathbf F\cdot\nabla)\boldsymbol v(y+s\mathbf F)$
in the distributional sense, hence
\[
(\mathbf F\cdot\nabla)\boldsymbol v(\boldsymbol x)={{} \boldsymbol h(\boldsymbol x).}
\]
Moreover, $\boldsymbol v\in L^2_{\mathrm{loc}}(\mathbb R^d;\mathbb R^d)$. Indeed, for any bounded set $K\subset\mathbb R^d$,
the set $\Phi^{-1}(K)\subset S\times\mathbb R$ is bounded, so there is a constant $T_K>0$ such that $|\tau(\boldsymbol x)|\le T_K$ for $\boldsymbol x\in K$; then
Cauchy--Schwarz in $s$ yields
\[
|\boldsymbol v(\boldsymbol x)|^2
\le 2T_K \int_{-T_K}^{T_K} |\boldsymbol h(\pi(\boldsymbol x)+s\mathbf F)|^2\,ds,
\]
and integrating over $K$ and changing variables $\boldsymbol x=\Phi(y,s)$ shows $\int_K |\boldsymbol v|^2<\infty$
because $\boldsymbol h\in L^2(\mathbb R^d)$ and $\Phi^{-1}(K)$ has finite measure.

The same construction gives the stated representation for arbitrary trace data.
Given any $\boldsymbol v_S\in L^2(S;\mathbb R^d)$, define
\begin{equation}\label{eq:v_def_Rd_vs}
\boldsymbol v(\boldsymbol x):=\boldsymbol v_S(\pi(\boldsymbol x)) +\int_0^{\tau(\boldsymbol x)} \boldsymbol h(\pi(\boldsymbol x)+s\mathbf F)\,ds.
\end{equation}
Then $\boldsymbol v|_S=\boldsymbol v_S$ since $\tau(\boldsymbol x)=0$ for $\boldsymbol x\in S$, and repeating the preceding argument along each
characteristic shows that $\boldsymbol v$ satisfies $(\mathbf F\cdot\nabla)\boldsymbol v={{} \boldsymbol h}$
in $\mathcal D'(\mathbb R^d)$ and belongs to $L^2_{\mathrm{loc}}(\mathbb{R}^d; \mathbb{R}^d)$.

We now prove part \textit{(ii)}. Let $\boldsymbol v_1,\boldsymbol v_2\in L^2_{\mathrm{loc}}(\mathbb R^d;\mathbb R^d)$ satisfy
$(\mathbf F\cdot\nabla)\boldsymbol v_j={{} \boldsymbol h}$ in $\mathcal D'(\mathbb R^d)$ and
$\boldsymbol v_1|_S=\boldsymbol v_2|_S=\boldsymbol v_S$. Set $\boldsymbol w:=\boldsymbol v_1-\boldsymbol v_2$.
Then $(\mathbf F\cdot\nabla)\boldsymbol w=0$ in $\mathcal D'(\mathbb R^d)$ and $\boldsymbol w|_S=0$.
For fixed $y\in S$, the function $s\mapsto \boldsymbol w(y+s\mathbf F)$ is constant in $s$
(since its distributional derivative is zero), and vanishes at $s=0$ because $\boldsymbol w|_S=0$.
Hence $\boldsymbol w(y+s\mathbf F)=0$ for all $s$ and for a.e.\ $y\in S$, and therefore $\boldsymbol w=0$
a.e.\ in $\mathbb R^d$ by the bijectivity of $\Phi$. This proves uniqueness for the prescribed trace.

Finally, we prove part \textit{(iii)}. Assume $\nabla\cdot\boldsymbol h=0$ in $\mathcal D'(\mathbb R^d)$. Taking divergence of
$(\mathbf F\cdot\nabla)\boldsymbol v={{} \boldsymbol h}$ yields
\[
(\mathbf F\cdot\nabla)(\nabla\cdot\boldsymbol v)={{} \nabla\cdot\boldsymbol h}=0\,.
\]
Here, we used the fact that since $\mathbf F$ is constant, the divergence operator commutes with $\mathbf F\cdot\nabla$. As a result, $\nabla\cdot\boldsymbol v$ is constant along characteristics. Since it vanishes on $S$, it vanishes identically in $\mathbb R^d$, and hence
$\nabla\cdot\boldsymbol v=0$ in $\mathcal D'(\mathbb R^d)$.
\end{proof}

\end{document}
